
\begin{multicols}{2}

Many European languages run the risk of becoming victims of the digital age because they are underrepresented and under-resourced online. Huge regional market opportunities remain untapped today because of language barriers. If we do not take action now, many European citizens will become socially and economically disadvantaged because they speak their native language.

Innovative language technology (LT) is an intermediary that will enable European citizens to participate in an egalitarian, inclusive and economically successful knowledge and information society. Multilingual language technology will be a gateway for instantaneous, cheap and effortless communication and interaction across language boundaries.

Today, language services are primarily offered by commercial providers from the US. Google Translate, a free service, is just one example. The recent success of Watson, an IBM computer system that won an episode of the Jeopardy game show against human candidates, illustrates the immense potential of language technology. As Europeans, we have to ask ourselves several urgent questions:

\begin{itemize}
\item Should our communications and knowledge infrastructure be dependent upon monopolistic companies?
\item Can we truly rely on language-related services that can be immediately switched off by others?
\item Are third parties from other continents willing to address our translation problems and other issues that relate to European multilingualism?
\item Can our European cultural background help shape the knowledge society by offering better, more secure, more precise, more innovative and more robust high-quality technology?
\end{itemize}

This whitepaper for the Bulgarian language demonstrates that a language technology research and development community exists in Bulgaria. However, many large companies have stopped their activities in language technology research and development, which is nowadays almost exclusively performed by small and medium enterprises that can hardly compete on the global market. 

Although a number of technologies and resources for Standard Bulgarian exist, there are fewer technologies and resources for the Bulgarian language than for the English language. The existing technologies and resources are also poorer in quality. 

According to the assessment detailed in this report, immediate action must occur before any breakthroughs for the Bulgarian language can be achieved.
\end{multicols}

\clearpage

\ssection[Risk for Our Languages and a Challenge for Language Technology]{Risk for Our Languages and a Challenge for Language Technology}

\begin{multicols}{2}

We are witnesses to a digital revolution that is dramatically impacting communication and society. Recent developments in digital information and communication technology are sometimes compared to Gutenberg’s invention of the printing press. What can this analogy tell us about the future of the European information society and our languages in particular?

After Gutenberg’s invention, real breakthroughs in communication and knowledge exchange were accomplished by efforts such as Luther’s translation of the Bible into vernacular language. In subsequent centuries, cultural techniques have been developed to better handle language processing and knowledge exchange:

\begin{itemize}
\item the orthographic and grammatical standardisation of major languages enabled the rapid dissemination of new scientific and intellectual ideas;
\item the development of official languages made it possible for citizens to communicate within certain (often political) boundaries;
\item the teaching and translation of languages enabled exchanges across languages;
\item the creation of editorial and bibliographic guidelines assured the quality and availability of printed material;
\item the creation of different media like newspapers, radio, television, books, and other formats satisfied different communication needs. 
\end{itemize}

In the past twenty years, information technology has helped to automate and facilitate many of the processes:

\begin{itemize}
\item desktop publishing software has replaced typewriting and typesetting;
\item Microsoft PowerPoint has replaced overhead projector transparencies;
\item e-mail send and receive documents faster than a fax machine;
\item Skype offers cheap Internet phone calls and hosts virtual meetings;
\item audio and video encoding formats make it easy to exchange multimedia content;
\item search engines provide keyword-based access to web pages;
\item online services like Google Translate produce quick, approximate translations;
\item social media platforms such as Facebook, Twitter, and Google+ facilitate communication, collaboration, and information sharing.
\end{itemize}

Although such tools and applications are helpful, they are not yet capable of supporting a sustainable, multilingual European society for all where information and goods can flow freely. 

\subsection[Language Borders Hold back the European Information Society]{Language Borders\newline Hold back the European Information Society}

We cannot predict exactly what the future information society will look like. But there is a strong likelihood that the revolution in communication technology is bringing people speaking different languages together in new ways. This is putting pressure on individuals to learn new languages and especially on developers to create new technology applications to ensure mutual understanding and access to shareable knowledge. In a global economic and information space, more languages, speakers and content interact more quickly with new types of media. The current popularity of social media (Wikipedia, Facebook, Twitter, YouTube, and, recently, Google+) is only the tip of the iceberg.

Today, we can transmit gigabytes of text around the world in a few seconds before we recognise that it is in a language we do not understand. According to a recent report from the European Commission, 57\% of Internet users in Europe purchase goods and services in languages that are not their native language. (English is the most common foreign language followed by French, German and Spanish.) 55\% of users read content in a foreign language while only 35\% use another language to write e-mails or post comments on the Web \cite{EC1}. A few years ago, English might have been the lingua franca of the Web—the vast majority of content on the Web was in English—but the situation has now drastically changed. The amount of online content in other European (as well as Asian and Middle Eastern) languages has exploded.

Surprisingly, this ubiquitous digital divide due to language borders has not gained much public attention; yet, it raises a very pressing question: Which European languages will thrive in the networked information and knowledge society, and which are doomed to disappear?

\subsection{Our Languages at Risk}

While the printing press helped step up the exchange of information in Europe, it also led to the extinction of many European languages. Regional and minority languages were rarely printed and languages such as Cornish and Dalmatian were limited to oral forms of transmission, which in turn restricted their scope of use. Will the Internet have the same impact on our languages?

Europe’s approximately 80 languages are one of its richest and most important cultural assets, and a vital part of its unique social model \cite{EC2}. While languages such as English and Spanish are likely to survive in the emerging digital marketplace, many European languages could become irrelevant in a networked society. This would weaken Europe’s global standing, and run counter to the strategic goal of ensuring equal participation for every European citizen regardless of language. According to a UNESCO report on multilingualism, languages are an essential medium for the enjoyment of fundamental rights, such as political expression, education and participation in society \cite{Unesco1}.

\subsection{Language Technology is a Key Enabling Technology}

In the past, investment efforts in language preservation focused on language education and translation. According to one estimate, the European market for translation, interpretation, software localisation and website globalisation was €8.4 billion in 2008 and is expected to grow by 10\% per annum \cite{EC3}. Yet this figure covers just a small proportion of current and future needs in communicating between languages. The most compelling solution for ensuring the breadth and depth of language usage in Europe tomorrow is to use appropriate technology, just as we use technology to solve our transport, energy and disability needs among others.

Digital language technology (targeting all forms of written text and spoken discourse) helps people collaborate, conduct business, share knowledge and participate in social and political debate regardless of language barriers and computer skills. It often operates invisibly inside complex software systems to help us:
\begin{itemize}
\item find information with an Internet search engine;
\item check spelling and grammar in a word processor;
\item view product recommendations in an online shop;
\item hear the verbal instructions of a car navigation system;
\item translate web pages via an online service.
\end{itemize}

Language technology consists of a number of core applications that enable processes within a larger application framework. The purpose of the META-NET language white papers is to focus on how ready these core technologies are for each European language. 

To maintain our position in the frontline of global innovation, Europe will need language technology adapted to all European languages that is robust, affordable and tightly integrated within key software environments. Without language technology, we will not be able to achieve a really effective interactive, multimedia and multilingual user experience in the near future.


\subsection{Opportunities for Language Technology}


In the world of print, the technology breakthrough was the rapid duplication of an image of a text (a page) using a suitably powered printing press. Human beings had to do the hard work of looking up, reading, translating, and summarising knowledge. We had to wait until Edison to record spoken language – and again his technology simply made analogue copies.

Digital language technology can now automate the very processes of translation, content production, and knowledge management for all European languages. It can also empower intuitive language/speech-based interfaces for household electronics, machinery, vehicles, computers and robots. Real-world commercial and industrial applications are still in the early stages of development, yet R\&D achievements are creating a genuine window of opportunity. For example, machine translation is already reasonably accurate in specific domains, and experimental applications provide multilingual information and knowledge management as well as content production in many European languages. 

As with most technologies, the first language applications such as voice-based user interfaces and dialogue systems were developed for highly specialised domains, and often exhibit limited performance. But there are huge market opportunities in the education and entertainment industries for integrating language technologies into games, cultural heritage sites, edutainment packages, libraries, simulation environments and training programmes. Mobile information services, computer-assisted language learning software, eLearning environments, self-assessment tools and plagiarism detection software are just some of the application areas where language technology can play an important role. The popularity of social media applications like Twitter and Facebook suggest a further need for sophisticated language technologies that can monitor posts, summarise discussions, suggest opinion trends, detect emotional responses, identify copyright infringements or track misuse.

Language technology represents a tremendous opportunity for the European Union. It can help address the complex issue of multilingualism in Europe – the fact that different languages coexist naturally in European businesses, organisations and schools. But citizens need to communicate across these language borders criss-crossing the European Common Market, and language technology can help overcome this final barrier while supporting the free and open use of individual languages. Looking even further forward, innovative European multilingual language technology will provide a benchmark for our global partners when they begin to enable their own multilingual communities. Language technology can be seen as a form of ‘assistive’ technology that helps overcome the ‘disability’ of linguistic diversity and make language communities more accessible to each other.

Finally, one active field of research is the use of language technology for rescue operations in disaster areas, where performance can be a matter of life and death: Future intelligent robots with cross-lingual language capabilities have the potential to save lives. 


\subsection{Challenges Facing Language Technology}

Although language technology has made considerable progress in the last few years, the current pace of technological progress and product innovation is too slow. Widely-used technologies such as the spelling and grammar correctors in word processors are typically monolingual, and are only available for a handful of languages. Online machine translation services, although useful for quickly generating a reasonable approximation of a document’s contents, are fraught with difficulties when highly accurate and complete translations are required. Due to the complexity of human language, modelling our tongues in software and testing them in the real world is a long, costly business that requires sustained funding commitments. Europe must therefore maintain its pioneering role in facing the technology challenges of a multiple-language community by inventing new methods to accelerate development right across the map. These could include both computational advances and techniques su
 ch as crowdsourcing.

\subsection{Language Acquisition in Humans and Machines}

To illustrate how computers handle language and why it is difficult to program them to use it, let’s look briefly at the way humans acquire first and second languages, and then see how language technology systems work. 

Humans acquire language skills in two different ways. Babies acquire a language by listening to the real interactions between its parents, siblings and other family members. From the age of about two, children produce their first words and short phrases. This is only possible because humans have a genetic disposition to imitate and then rationalise what they hear. 

Learning a second language at an older age requires more effort, largely because the child is not immersed in a language community of native speakers. At school, foreign languages are usually acquired by learning grammatical structure, vocabulary and spelling using drills that describe linguistic knowledge in terms of abstract rules, tables and examples. Learning a foreign language gets harder with age.

The two main types of language technology systems ‘acquire’ language capabilities in a similar manner. Statistical (or ‘data-driven’) approaches obtain linguistic knowledge from vast collections of concrete example texts. While it is sufficient to use text in a single language for training, e.g., a spell checker, parallel texts in two (or more) languages have to be available for training a machine translation system. The machine learning algorithm then "learns" patterns of how words, short phrases and complete sentences are translated. 

This statistical approach can require millions of sentences and performance quality increases with the amount of text analysed. This is one reason why search engine providers are eager to collect as much written material as possible. Spelling correction in word processors, and services such as Google Search and Google Translate all rely on statistical approaches. The great advantage of statistics is that the machine learns fast in continuous series of training cycles, even though quality can vary arbitrarily.

The second approach to language technology and machine translation in particular is to build rule-based systems. Experts in the fields of linguistics, computational linguistics and computer science first have to encode grammatical analyses (translation rules) and compile vocabulary lists (lexicons). This is very time consuming and labour intensive. Some of the leading rule-based machine translation systems have been under constant development for more than twenty years. The great advantage of rule-based systems is that the experts have more detailed control over the language processing. This makes it possible to systematically correct mistakes in the software and give detailed feedback to the user, especially when rule-based systems are used for language learning. But due to the high cost of this work, rule-based language technology has so far only been developed for major languages. 

As the strengths and weaknesses of statistical and rule-based systems tend to be complementary, current research focuses on hybrid approaches that combine the two methodologies. However, these approaches have so far been less successful in industrial applications than in the research lab. 

As we have seen in this chapter, many applications widely used in today’s information society rely heavily on language technology. Due to its multilingual community, this is particularly true of Europe’s economic and information space. Although language technology has made considerable progress in the last few years, there is still huge potential in improving the quality of language technology systems. In the following, we will describe the role of Bulgarian in European information society and assess the current state of language technology for the Bulgarian language. 

\end{multicols}

\clearpage

\ssection[Bulgarian in the European Information Society]{Bulgarian in the European Information Society}

\begin{multicols}{2}

\subsection{General Facts}

Bulgarian is the official language of the Republic of Bulgaria. 

Bulgarian is spoken by approximately 9 million native speakers, mainly in Bulgaria. but also in Greece, Macedonia, Romania, Serbia, Turkey (Europe), Ukraine, Australia, Canada, USA, Germany and Spain. It is also spoken in Croatia, Czech Republic, Hungary, Israel, Moldova, Romania, Russian Federation (Europe), Slovakia \cite{Ethno}. Bulgarian communities in the neighbouring countries and communities of immigrants around the world are supported by the Government agency for Bulgarians Abroad \cite{ABA} at ministerial level.

Preliminary data provided by the National Statistical Institute \cite{NSI} from the population census of Bulgaria indicates that as of the 1st of February, 2011, the population was 7 351 633. The number of Bulgarians living in countries of the European Union will become clear in 2012 when the data from EU populations censuses is available.

The official alphabet is Cyrillic. Bulgarian is the first Slavic language with its own writing system which dates from the 9th century. In 886 AD the Glagolitic alphabet was introduced in Bulgaria. Glagolitic was created by Sts. Cyril and Methodius but was gradually replaced with the Cyrillic alphabet, created at the literary schools of Ohrid and Preslav in the beginning of the 10th century. On the 1st of January 2007 when Bulgaria became a full member of the European Union, Cyrillic became the third official alphabet of the European Union, following the Latin and Greek alphabets.
 
Bulgarian dialects are those regional, colourful varieties of the Bulgarian language spoken both inside and outside the borders of the country. The main isogloss separating the Bulgarian dialects into Eastern and Western is the "Yat Border" which marks the different mutations of the Old Bulgarian "yat" form (\includegraphics[scale=0.4]{../_media/bulgarian/yat}),  pronounced in certain conditions as either /’a/ or /e/ to the east ({\it /b’al/}, but plural {\it /beli/}, [white]) and strictly as /e/ to the west of it ({\it /bel/}, plural {\it /beli/}). 

\subsection{Particularities of the Bulgarian Language}

Bulgarian belongs to the family of South Slavic languages and forms part of the Balkan linguistic union (Balkan Sprachbund). Consequently Bulgarian displays similarities to both language groups. As a Slavic language Bulgarian possesses a rich inflectional and derivational morphology, verb aspect pairs, etc. However, due to the mutual influence of Balkan languages Bulgarian has lost noun cases (except vocative) and completely has lost the infinitive form.

In addition to a common lexicon some of the most important grammatical features of Bulgarian, which distinguish it as a Slavic language, are (some of these features are also shared with other Balkan languages):
\begin{itemize}
\item Rich system of inflections. The rich inflectional system poses specific difficulties to LT systems: lemmatisation, for example, might face the notorious problem of recognition of certain inflectional types that can belong to different parts of speech. Such a case of homography, for instance is the word \textit{{beli}} which can be: 
\begin{itemize}
\item plural of the noun \textit{{belya}} [trouble];
\item plural of the adjective \textit{{byal}} [white];
\item 3rd person singular presence of the verb \textit{{belya}} [peel];
\item 2nd or 3rd person singular aorist of the verb \textit{{belya}} [peel];
\item 2nd person singular imperative of the verb \textit{{belya}} [peel]. 
\end{itemize}

\item Rich derivational system: diminutive and augmentative nouns: \textit{{stol — stolche}} [chair — small seat]; feminine suffixes for nouns relating to persons (professions): \textit{{uchitel — uchitelka}} [teacher — female teacher]; etc.

\item Aspectual verb pairs: \textit{{rodya — razhdam}} [to give birth – to be giving birth].

\item A pro-drop language, i.e. subjects are typically not overtly expressed whenever they are inferable from context: \textit{{Az cheta kniga — Cheta kniga}} [I am reading a book]. 
\end{itemize}

During the historical development of Bulgarian and as a result of its contacts with neighbouring non-Slavic languages in the Balkans, significant changes have occurred when compared with the other Slavic languages. Typical features of the Balkan Sprachbund are:
\begin{itemize}
\item Simplification of the nominal case structure — only relics of the accusative and dative cases remain in the system of pronouns and the vocative form in certain masculine and feminine animate nouns or persons.

\item Existence of post-fixed definite article in Bulgarian (\textit{{zhena — zhenata}} [a woman — the woman]), Albanian, Aroumanian and Romanian. This phenomenon arose most probably under the influence of Old Bulgarian since it does not exist in Modern Greek, where the definite article is prefixed.

\item Loss of the infinitive and its replacement with finite verb forms, etc. 
\end{itemize}

Bulgarian orthography is far more transparent than, for example, English orthography. The pronunciation of a word can be easily divined from its written form. Nonetheless, a number of the writing system's rules combined with pronunciation rules (assimilation between voiced and voiceless consonants, reduction of vowels, etc.) lead to potential homophony: \textit{{kos} /kos/} [blackbird] and \textit{{koz /kos/}} [trump]. In this way, a number of spellings could represent the same pronunciation. This is a source of potential errors that spell and grammar checkers are able to detect and correct.

Capitalisation in Bulgarian is sparse compared with other languages. In general, only personal and place names, some abbreviations (e.g. \textit{Stara Zagora}), the first word (only) in the title of a book, movie, song, etc., and the first word in a sentence are capitalised, as are names of companies, government bodies, etc. Names of nationalities or languages are not capitalised, nor are days of the week and months of the year.
 
Bulgarian exhibits a number of specific characteristics that contribute to the richness of the language but can also be a challenge for the computational processing of natural language. For example, it has a relatively free word order – while the order of adjectives as pre-positive modifiers of nouns or of adverbs as pre-positive modifiers of adjectives and other adverbs is fixed, the order of the subject and object of the verb is relatively loose. Thus in a sentence of three constituents (subject, verb, object), all six permutations are possible, in a sentence containing four constituents (subject, verb, direct object, indirect object) — twenty four permutations, and so on. For example, in the following sentence:

 \textit{{Divata kotka goni zloto kuche tsyala sutrin}} [The wild cat chases the bad dog the whole morning] it is not clear without a broader context which constituent is the subject — \textit{{divata kotka}} [wild-the cat] or \textit{{zloto kuche}} [bad-the dog]. Unlike other languages which show a relatively loose word order (the other Slavonic languages, for example) Bulgarian does not possess nominal case inflection to indicate the syntactic relations between the words.

Yet another characteristic feature of Bulgarian which poses difficulty for syntactic parsing is the free omission of the subject which, when combined with the possibility of shifting the positions of the subject and object, makes the task even harder. For example the following sentence: \textit{{Izvika Anna}} [*called Anna] could mean that Anna called somebody, i.e.  \textit{Anna} is the subject, or that Anna had been called by someone else, i.e. that  \textit{Anna} is the object (there is no morphological marker to distinguish the subject from the object, except in the case of the graphical representation of the definite article in masculine nouns in the subject position: \textit{{popita uchitelyat}} [*asked teacher] means that the teacher asked something, while \textit{{popita uchitelya}} [*asked teacher] means that the teacher was asked something. 

The rather flexible word order which when combined with the lack of morphological distinction for nominal cases and subject omission is a real challenge for natural language processing of Bulgarian.

\subsection{Recent Developments}

Since 1990 foreign (mostly American) movies and television series have begun to dominate Bulgarian broadcasting. Foreign films and series are either dubbed (mainly by the national television and other bigger television companies, nationally broadcast) or subtitled (mainly at smaller private television companies). However, whatever the way in which they are translated, the strong presence of foreign way of life in the media has influenced Bulgarian culture and language. Due to the continuing triumph of English and American music since the 1960s (much more since 1990), Bulgarians were exposed to a lot of English during their adolescence. English quickly acquired the status of a ‘cool/hip’ language in certain areas, which it has kept up to the present day. 

During the last 20 years there has been a noticeable trend towards the "internationalisation" of the Bulgarian lexicon as a result of the influence of English. Bulgarian has accepted new words, meanings and collocations originating predominantly from English (so-called Anglicisms), although many others have been taken from other European and non-European languages.

In the Dictionary of New Words in Bulgarian (2010)  \cite{NWDict} which records new words, phraseology and terminological expressions from the past 20 years, about 4300 new lexical units have been registered. Of these about a quarter (about 1020) are borrowings from English. There are a number of terminological areas in which the lexicon has developed almost entirely under the influence of English: computer technology and the internet (\textit{{fayl}} [file], \textit{{sayt}} [site]), finances, economics and business (\textit{{dilar}} [dealer], \textit{{broker}} [broker]), contemporary music (\textit{{didzhey}} [dj], \textit{{tehno}} [techno], \textit{{klip}} [clip]), sport (\textit{{dzhoging}} [jogging], \textit{{bodibilding}} [bodybuilding]). The influx of English borrowings has also been seen in commonly used words — e.g. \textit{{toster}} [toaster], \textit{{stiker}} [sticker], \textit{{bodigard}} [bodyguard], and teenage slang.

The dictionary also records over 700 new word meanings, the majority of which have arisen under the influence of English. These are semantic calques, many of which are in the area of computer terminology: \textit{{mishka}} [mouse], \textit{{papka}} [folder], \textit{{glasova poshta}} [voice mail], etc. The influence of English in these cases is not always obvious since it typically affects the word senses, rather then formal composition of words. 

Many of the new borrowings from English cause difficulties for speakers of Bulgarian. Some are difficult to pronounce – \textit{{blokbastar}} [blockbuster], \textit{{marchandayzing}} [merchandising], while others are difficult to adapt morphologically — some words give rise to uncertainty when used in the plural for example, \textit{{bodigardi}} or \textit{{bodigardove}} [bodyguards], \textit{{chipseti}} or \textit{{chipsetove}} [chipsets]. Older generation Bulgarians who do not speak English find these borrowings hard to understand. 

A new phenomenon unknown in Bulgarian before the 1990’s is the graphic representation of many foreign words in the Latin alphabet, in particular English borrowings. A particular case in point are abbreviations (such as CV, CD-ROM, PR), as well as widely used words such as \textit{internet}, or even prefixed component parts such \textit{e-} for electronic. The use of Latin is particularly prevalent in certain specialised areas such as: computer technology, in the names of companies, commercial establishments etc., but also in commonly used language, such as advertising. Some people are worried that the use of the Latin alphabet for writing Bulgarian (called \textit{{shlyokavitsa} /shl’okavitsa/} and used mainly in unofficial modes — sms, e-mails, etc.) \cite{shlyokavica} will somehow affect the quality of spoken and written Bulgarian and eliminate the use of the distinctive Cyrillic alphabet.

The few examples given demonstrate the importance of raising awareness of developments which might lead to the risk of excluding large groups from participation in the information society, namely those who are not familiar with English.

\subsection{Language Cultivation in Bulgaria}

Bulgarian is the official language in the Republic of Bulgaria, as stated in the Bulgarian constitution. Constituted by decree of the Council of Ministers, the Institute for Bulgarian Language of the Bulgarian Academic of Science is the official institution which observes changes in the Bulgarian language, determines literary norms and reflects these changes in both orthography and speech. Its primary tasks include research in Bulgarian linguistics, general, theoretical, applied and computational linguistics, as well as the preparation of a comprehensive dictionary of the Bulgarian language, and the maintenance of its archival materials. Other research projects investigate Bulgarian dialects in and outside Bulgaria, including issues of language policy within the framework of European integration. Further tasks include the assembly of linguistic corpora and databases, and laying the linguistic groundwork for computational software and applications. Over the past 60 years as a 
 product of these functions the Institute has published the academic journal "Balgarski ezik" [Bulgarian Language] and has provided information and advice through its Language service. Regular broadcasts on Bulgarian National Radio are also provided and until recently and collaboration with Bulgarian National Television has produced "Ezik moy" [My Language] broadcast. 

In parallel with these activities many of the national and local daily newspapers carry articles concerning matters of language: for example the regular column in the "Trud" national daily newspaper, written by professors of the University of Sofia. Since 2002 the University of Sofia has published "Rodna rech" [Native tongue], a journal whose main aim is to respond to a range of current language issues and assist in raising linguistic culture amongst Bulgarian society. 

The first academic spelling dictionary was published in 1983 and reflected the accepted standards of spelling and speech. In 2002 a new spelling dictionary was published which was augmented by new words and certain changes which mirrored the development of the language. As well as dictionaries which serve to define standard language, other spelling and pronunciation dictionaries are also published to reflect standard usage. Numerous periodicals connected with problems of spelling and pronunciation are also published on a regular basis. 

Most of these activities are of purely academic interest and are not sufficiently popular among the younger population in particular. Media language seeks to attract and entertain and in many cases deviates from the proscriptions of standard usage. No official quotas regulate the percentage of Bulgarian language music song (even on the National Radio and Television). This should be compared to the regulations for example in France, Hungary, Slovenia. Nevertheless, the current Law on Radio and Television states that Bulgarian National Radio shall set aside for the creation and performance of Bulgarian music and radio drama not less than 5\% of the subsidy of the state budget; while Bulgarian National Television shall set aside not less than 10\% of the same subsidy for Bulgarian television film production. 

The most recent orthographic reform was carried out in 1945. Its aim was to bring the spelling of Bulgarian up to date and consisted of the removal of certain letters related historical orthographic forms. On a number of occasions members of the Bulgarian National Assembly have submitted draft versions of a new Law on Bulgarian Language primarily in the aims of preserving the purity of the language. These have led to long and rather emotional debates. In 2007 the Law on Transliteration was passed in the aims of standardising the diverse practices of rendering words (personal names) in Latin letters.

The situation of the Bulgarian language is disadvantageous when compared to languages such as French, which is strongly promoted by the global community of French-speaking peoples within the so-called Francophone Union. The wide usage of Language technology can make an important contribution in this area by offering media, internet and mobile communications sophisticated language services — spell and grammar checkers, style correction, dictionary checks for synonyms, and the correct pronunciation of words.

\subsection{Language in Education}

From the 19th century onwards Bulgarian language and literature has had a very important role in the education. According to the legislation in Bulgaria, all education and teaching provided as part of the current state curriculum, from pre-school through to university level, must be in Bulgarian. Special arrangements exist for children whose mother tongue is not Bulgarian. The study of Bulgarian is compulsory for the elementary and secondary school.

In PISA‘2009 \cite{oecd} (the capacity to use scientific knowledge, to identify questions and to draw evidence) Bulgaria is in 46th place out of 65 countries. In PISA‘2006 Bulgaria is in 43th place out of 55 countries. In PISA‘2000 — in 33th place out of 41 countries. 
In PIRLS‘2006 \cite{nces} (reading literacy as the ability to understand and use those written language forms required by society) Bulgaria was in 14th place with 547 points. Again the results were lower than in PIRLS’2001 — 4th place with 550 points.

The Bulgarian PISA and PIRLS results can be used as an indicator (or International Benchmark) to determine to what extent international educational standards are satisfied within the National School Curriculum. In 2009 the Bulgarian students demonstrated basic literacy and almost half of them found it difficult to interpret and analyse a text. 

The insufficiencies of Bulgarian language teaching in high schools (for example) can be summarised in a number of points: insufficient allocated time – 36 hours (lessons) annually; non-communicative organisation of the teaching process; inadequate content.

According to the most recent State Educational Requirements Bulgarian language teaching is conducted within the framework of a cultural and education study sector — Bulgarian Language and Literature. This sector is traditional within Bulgarian schools and the universities train specialist — middle and high-school teachers in this subject. One of the ways of increasing the effectiveness of Bulgarian language teaching is for it to be focused upon as a specific and important scientific area. Although traditionally seen as a humanitarian discipline, linguistics is concerned with the formulation of rules according to which language units are inter-combined, and is thus close to the sciences. 

At university level, there is a general shortage of courses in Bulgarian (at some Universities) that would enable future experts for successful professional communication and appropriate functional literacy.

Language skills are the key qualification needed in education as well as for personal and professional communication. Increasing the volume of Bulgarian language teaching in schools is one possible step towards providing students with the language skills required for active participation in society. Language technology can make an important contribution here by offering so-called computer-assisted language learning (CALL) systems. Such systems allow students to experience language through play; for example by linking special vocabulary in an electronic text to comprehensible definitions or to audio or video files supplying additional information, e.g., the pronunciation of a word. 

\subsection{International Aspects}

The usage and influence of the Bulgarian language beyond the borders of Bulgaria and its speakers is limited. For many years the Ministry of Education, Youth and Science has held competitions to send Bulgarian university lecturers to work as lectors in Bulgarian language, literature and culture in a number of foreign universities where Bulgarian is studied. The Ministry of Culture established the National Culture Fund which holds regular competitions for the translation of Bulgarian literature into foreign languages. Translations of the works of classical Bulgarian writers, such as Hristo Botev and Ivan Vazov, into almost all European languages have, of course, been in existence for many years. 

Bulgaria has a great number of world famous singers (Nikolay Gyaurov, Gena Dimitrova, Boris Hristov, Valya Balkanska), authors and artists (Hristo Yavashev, Donyo Donev, Yuliya Krasteva, Tsvetan Todorov) among many others. The poet, Pencho Slaveikov (1866-1912), a central figure in Bulgarian literature, is the only Bulgarian to have been nominated for the Nobel Prize. Outside certain narrow specialised circles overseas where the Bulgarian languages is taught and studied, Bulgarian is an unknown and exotic language.

As everywhere in the scientific world, Bulgarian scientists face a great deal of pressure to publish in visible (usually international) journals, most of which are now in the English language. The situation is similar in the business world. In many large and internationally active companies, English has become the lingua franca, both in written and oral communication. At the same time 44\% of mature Bulgarians do not speak a foreign languages according to research published by the European statistical service, Eurostat \cite{epp}, in 2009. 

Bulgarian has acquired the status as an official administrative language of the European Union on the same basis as English, German and French, since the European Union is based on solidarity and equality amongst its members. Since the 1st January 2007, Bulgarian has been used in the following situations in the context of relations between Bulgaria and the EU:

\begin{itemize}
\item The official bulletin containing the rights of citizens and the texts of EU law is published in Bulgarian.
\item The Bulgarian authorities are entitled to speak in Bulgarian in Council of the European Union.
\item Bulgarian citizens are entitled to use Bulgarian in their correspondence with the European institutions.
\end{itemize}

The fact of Bulgaria’s membership of the European Union together with the idea of unity and diversity, globalisation while preserving national identity, provides a real opportunity for the egalitarian use of Bulgarian together with the major European languages.

Language technology can address this challenge from a different perspective by offering services such as machine translation or cross-lingual information retrieval to foreign language text and thus help diminish personal and economic disadvantages naturally faced by non-native speakers of English.

\subsection{Bulgarian on the Internet}

Bulgarian internet usership in 2009 increased by 31\% in comparison with 2007 and already 46\% of the total population uses the internet. According to a study by gemiusAudience \cite{gemius}, published in the report "Do you CEE?" \cite{inetcee} Bulgaria is amongst the countries with the highest percentage of internet penetration.

According to data published by internetworldstats.co \cite{inetworldstat} there are about 3.5 million internet users in Bulgaria, and after the statistics published by Gemius the growth of sites observed by analysts is almost 10.7\% on an annual basis. In 2010 there was a further 5\% increase in usership.

In addition to the ubiquitous international web sites, the most popular web sites on the Bulgarian part of the Internet are Bulgarian news portals (dir.bg, gbg.bg news.bg, etc.). Bulgarian Wikipedia as an important source for natural language processing contains app. 117\,000 articles, a considerably smaller size than the biggest Wikipedias – English, German and French – but in the number of articles it is in the 34th position \cite{metadata} among 270 Wikipedias in other languages. 

It is often claimed that English dominates computers and the internet, and that those wishing to use either must first learn English. That may have been true in the early days of the technology but lack of English is no longer the barrier it once was. What began as an anglophone phenomenon has rapidly become a multilingual affair. Software has been made capable of displaying many different kinds of script. Many corporate websites now employ multilingual strategies making choice of language a ‘user preference’. Machine translation of web content is only a mouse-click away.

For language technology, the growing importance of the internet is important in two ways. On the one hand, the large amount of digitally available language data represents a rich source for analysing the usage of natural language, in particular by collecting statistical information. On the other hand, the internet offers a wide range of application areas involving language technology. 

One important aspect of equal opportunities the Law on Equal Opportunities for the Disabled, which came into force in 2002, and addresses the issue of barrier-free information technology. It enjoins public agencies to make sure that the disabled can use their websites and Internet services without any restrictions. User-friendly language technology tools are a key solution to this requirement by offering for example speech synthesis to enunciate the content of web pages for the blind.

Internet users and providers of web content can also use language technology in less obvious ways, for example, by automatically translating web page contents from one language into another. Despite the high cost of manually translating this content, comparatively little language technology has been developed and applied to the issue of website translation in light of the supposed need. This may be due to the complexity of the Bulgarian language and to the range of different technologies involved in typical applications.

The next chapter gives an introduction to language technology and its core application areas, together with an evaluation of current language technology support for Bulgarian. 
\end{multicols}

\clearpage

\ssection[Language Technology Support for Bulgarian]{Language Technology Support\newline for Bulgarian}

\begin{multicols}{2}

Language technologies are software systems designed to handle human language and are therefore often called "human language technology". Human language comes in spoken and written forms. While speech is the oldest and in terms of human evolution the most natural form of language communication, complex information and most human knowledge is stored and transmitted in written texts. Speech and text technologies process or produce these different forms of language, though they both use dictionaries and rules of grammar and semantics. This means that language technology (LT) links language to various forms of knowledge, independently of the media (speech or text) it is expressed in. The figure illustrates the LT landscape. When we communicate, we combine language with other modes of communication and information media – for example speaking can involve gestures and facial expressions. Digital texts link to pictures and sounds. Movies may contain language in spoken and written f
 orm. In other words, speech and text technologies overlap and interact with other technologies that facilitate processing of multimodal communication and multimedia documents. 

In the following, we will discuss the main application areas of language technology, i.e., language checking, web search, speech technology, and machine translation. This includes applications and basic technologies such as 
\begin{itemize}
\item spelling correction;
\item authoring support;
\item computer-assisted language learning;
\item information retrieval;
\item information extraction;
\item text summarisation;
\item question answering;
\item speech recognition; 
\item speech synthesis.
\end{itemize}

Before discussing the above application areas, we will shortly describe the architecture of a typical LT system. 

\subsection{Application Architectures}

Software applications for language processing typically consist of several components that mirror different aspects of language. The figure shows a highly simplified architecture that can be found in a typical text processing system. The first three modules handle the structure and meaning of the text input:
\begin{enumerate}
\item Pre-processing: cleans the data, analyses or removes formatting, detects the input language, and so on.
\item Grammatical analysis: finds the verb, its objects, modifiers and other parts of speech as well as detects the sentence structure.
\item Semantic analysis: performs disambiguation (i.e., computes the appropriate meaning of words in a given context); resolves anaphora (i.e., which pronouns refer to which nouns in the sentence) and substitute expressions; and represents the meaning of the sentence in a machine-readable way.
\end{enumerate}

After analysing the text, task-specific modules can perform other operations, such as automatic summarisation and database look-ups. This is a simplified and idealised description of the application architecture and illustrates the complexity of LT applications. 

After introducing the core application areas for language technology, we shall provide a brief overview of the state of LT research and education today, and end with an overview of past and present research programmes. We shall then present an expert estimate of core LT tools and resources in terms of various dimensions such as availability, maturity and quality. The general situation of LT for the Bulgarian language is summarised in a table.

\subsection{Core Application Areas}

In this section, we focus on the most important LT tools and resources, and give an overview of LT activities in Bulgaria. Tools and resources that are \uline{underlined} in the text can also be found in the table at the end of this chapter.  

\subsubsection{Language checking}

Anyone who has used a word processor such as Microsoft Word knows that it has a spelling checker that highlights spelling mistakes and proposes corrections. The first spelling correction programs compared a list of extracted words against a dictionary of correctly spelled words. Today these programs are far more sophisticated. Using language-dependent algorithms for \uline{text analysis}, they detect errors related to morphology (e.g., plural formation) as well as syntax–related errors, such as a missing verb or a conflict of verb-subject agreement (e.g., \textit{Tya *napisahme pismoto} [she *write a letter]). But most spell checkers will not find any errors in the following text:  

{\it
\hspace{0.5cm}Eye have a spelling chequer,

\hspace{0.5cm}It came with my Pea Sea.

\hspace{0.5cm}It plane lee marks four my revue

\hspace{0.5cm}Miss Steaks I can knot sea. \cite{zar1}
}

For handling this type of error, \uline{context} analysis is needed in many cases, e.g., to decide if a word needs to be written in upper case, as in:

{\it
\hspace{0.5cm}{Tya zhivee v \textbf{Stara} Zagora.}
}

\hspace{0.5cm}[She lives in Stara Zagora.]

{\it
\hspace{0.5cm}{Tya e \textbf{stara} zhena.}
}

\hspace{0.5cm}[She is an old woman.]


This type of analysis either needs to draw on language-specific \uline{grammars} laboriously coded into the software by experts, or on a statistical language model. In this case, a model calculates the probability of a particular word as it occurs in a specific position (e.g., between the words that precede and follow it). A statistical language model can be automatically created by using a large amount of (correct) language data (called a \uline{text corpus}). Most of these two approaches have been developed around data from English. Neither approach can transfer easily to Bulgarian because the language has a flexible word orderand a richer inflection system. 

Language checking is not limited to word processors; it is also used in "authoring support systems", i.e. software environments in which manuals and other documentation are written to special standards for complex IT, healthcare, engineering and other products. Fearing customer complaints about incorrect use and damage claims resulting from poorly understood instructions, companies are increasingly focusing on the quality of technical documentation while targeting the international market (via translation or localisation) at the same time. Advances in natural language processing have led to the development of authoring support software, which helps the writer of technical documentation use vocabulary and sentence structures that are consistent with industry rules and (corporate) terminology restrictions.

Besides spell checkers and authoring support, language checking is also important in the field of computer-assisted language learning. And language checking applications also automatically correct search engine queries, as found in Google's \textit{Did you mean...} suggestions. 

\subsubsection{Web search}

Searching the Web, intranets or digital libraries is probably the most widely used yet largely underdeveloped language technology application today. The Google search engine, which started in 1998, now handles about 80\% of all search queries \cite{spi1}. Since 2004, the verb googeln even has an entry in the Duden dictionary. The Google search interface and results page display has not significantly changed since the first version. Yet in the current version, Google offers spelling correction for misspelled words and  has now incorporated basic semantic search capabilities that can improve search accuracy by analysing the meaning of terms in a search query context  \cite{pc1}. The Google success story shows that a large volume of available data and efficient indexing techniques can deliver satisfactory results for a statistically-based approach.  

For more sophisticated information requests, it is essential to integrate deeper linguistic knowledge for \uline{text interpretation}. Experiments using \uline{lexical resources} such as machine-readable thesauri or ontological language resources (e.g., WordNet for English or BulNet for Bulgarian) have demonstrated improvements in finding pages using synonyms of the original search terms, such as {\it atomna energiya} [atomic energy] and {\it yadrena energiya} [nuclear energy], or even more loosely related terms. 

The next generation of search engines will have to include much more sophisticated language technology, in particular in order to deal with search queries consisting of a question or other sentence type rather than a list of keywords. For the query, \textit{"Give me a list of all companies that were taken over by other companies in the last five years"}, the LT system needs to analyse the sentence syntactically and semantically as well as provide an index to quickly retrieve relevant documents. A satisfactory answer will require syntactic parsing to analyse the grammatical structure of the sentence and determine that the user wants companies that have been acquired, not companies that acquired other companies. For the expression \textit{last five years}, the system needs to determine the relevant years. And, the query needs to be matched against a huge amount of unstructured data to find the piece or pieces of relevant information the user wants. This is called "information r
 etrieval", and involves searching and ranking relevant documents. To generate a list of companies, the system also needs to recognise a particular string of words in a document as a company name, a process called "named entity recognition".

A more demanding challenge is matching a query in one language with documents in another language. Cross-lingual information retrieval involves automatically translating the query into all possible source languages and then translating the results back into the target language. 

Now that data is increasingly found in non-textual formats, there is a need for services that deliver multimedia information retrieval by searching images, audio files and video data. In the case of audio and video files, a speech recognition module must convert the speech content into text (or into a phonetic representation) that can then be matched against a user query.

Certain Bulgarian portals have crawler software similar to those used by global search engines designed to index sites included within its categories. For example, Dir.bg, one of the first and largest web portals in Bulgaria launched a standalone service – Diri.bg. Diri (in Bulgarian {дири}) is an old word for ‘search’. 


Open source based technologies like Lucene and SOL are often used by search-focused companies to provide the basic search infrastructure. Other search-based companies rely on international search technologies like, e.g., FAST or Exalead. Developmental focus for these companies lies in providing add-ons and advanced search engines for special-interest portals by exploiting topic-relevant semantics. Due to the still high demands in processing power, such search engines are only economically usable on relatively small text corpora. Processing time easily exceeds that of a common statistical search engine as, e.g., provided by Google by a magnitude of thousands. These search engines also have high demand in topic-specific domain modelling, making it not feasible to use these mechanisms on web scale.

\subsubsection{Speech Technology}

Speech interaction technology is used to create interfaces that enable users to interact in spoken language instead of a graphical display, keyboard and mouse. Today, these voice user interfaces (VUI) are used for partially or fully automated telephone services provided by companies to customers, employees or partners. Business domains that rely heavily on VUIs include banking, supply chain, public transportation, and telecommunications. Other uses of speech interaction technology include interfaces to car navigation systems and the use of spoken language as an alternative to the graphical or touch-screen interfaces in smartphones. 

Speech interaction comprises four technologies: 

\begin{enumerate}
\item Automatic \uline{speech recognition} (ASR) determines which words are actually spoken in a given sequence of sounds uttered by a user.
\item Natural language understanding analyses the syntactic structure of a user’s utterance and interprets it according to the system in question.
\item Dialogue management determines which action to take given the user input and system functionality.
\item \uline{Speech synthesis} (text-to-speech or TTS) transforms the system’s reply into sounds for the user. 
\end{enumerate}

One of the major challenges of ASR systems is to accurately recognise the words a user utters. This means restricting the range of possible user utterances to a limited set of keywords, or manually creating language models that cover a large range of natural language utterances. Using machine learning techniques, language models can also be generated automatically from \uline{speech corpora}, i.e. large collections of speech audio files and text transcriptions. Restricting utterances usually forces people to use the voice user interface in a rigid way and can damage user acceptance; but the creation, tuning and maintenance of rich language models will significantly increase costs. VUIs that employ language models and initially allow a user to express their intent more flexibly — prompted by a \textit{How may I help you?} greeting — tend to be automated and are better accepted by users. 

Companies tend to use pre-recorded utterances by professional speakers for generating the output of the voice user interface. For static utterances where the wording does not depend on particular contexts of use or personal user data, this can deliver a rich user experience. But more dynamic content in an utterance may suffer from unnatural intonation because bits of audio files have simply been strung together. Today’s TTS systems are getting better (though they can still be optimised) at producing natural-sounding dynamic utterances.  

Interfaces in the market for speech interaction technologies have been considerably standardised during the last decade in terms of their various technology components. There has also been strong market consolidation in speech recognition and speech synthesis. The national markets in the G20 countries (economically resilient countries with high populations) have been dominated by just five global players, with Nuance (USA) and Loquendo (Italy) being the most prominent players in Europe. In 2011, Nuance announced the acquisition of Loquendo, which represents a further step in market consolidation.

On the Bulgarian TTS market, there are a few Bulgarian text-to-speech systems. One of these is SpeechLab 2.0 provided free-of-charge to computer users with visual disabilities. Only a few companies such as "Ciela" – a Bulgarian publisher of legal literature — have developed their own system for Bulgarian speech recognition. Finally, within the area of Speech Interaction, the is as yet no real market for syntactic and semantic analysis-based core technologies. 

The demand for voice user interfaces in Bulgaria has grown fast in the last five years, driven by increasing demand for customer self-service, cost optimisation for automated telephone services, and the increasing acceptance of spoken language as a media for human-machine interaction. 

Looking forward, there will be significant changes due to the spread of smartphones as a new platform for managing customer relationships in addition to fixed telephones, the Internet and e-mail. This will also affect how speech interaction technology is used. In the long run, there will be fewer telephone-based VUIs and spoken language will play a far more central role as a user-friendly input for smartphones. This will be largely driven by stepped improvements in the accuracy of speaker-independent speech recognition via speech dictation services already offered as centralised services to smartphone users. 

\subsubsection{Machine Translation}

The idea of using digital computers to translate natural languages goes back to 1946 and was followed by substantial funding for research during the 1950s and again in the 1980s. Yet \uline{machine translation} (MT) still cannot meet its initial promise of across-the-board automated translation.  

The most basic approach to machine translation is to automatically replace the words in a text in one natural language by words in another language. This can be useful in subject domains that have a very restricted, formulaic language such as weather reports. But to produce a good translation of less standardised texts, larger text units (phrases, sentences, or even whole passages) need to be matched to their closest counterparts in the target language. The major difficulty is that human language is ambiguous. Ambiguity creates challenges on multiple levels, such as word sense disambiguation on the lexical level ({\it a jaguar} is a brand of car or an animal) or the assignment of case on the syntactic level, for example:

{\it
\hspace{0.5cm}{Politsayat nablyudava prestapnika s tele-skopa.}
}

\hspace{0.5cm}[The policeman observed the man with the telescope.]

{\it
\hspace{0.5cm}{Politsayat nablyudava prestapnika s pistoleta.}
}

\hspace{0.5cm}[The policeman observed the man with the revolver.]


One way to build an MT system is to use linguistic rules. For translations between closely related languages, a direct substitution translation may be feasible in cases like the above example. But, rule-based (or linguistic knowledge-driven) systems often analyse the input text and create an intermediary symbolic representation from which the text can be generated into the target language. The success of these methods is highly dependent on the availability of extensive lexicons with morphological, syntactic, and semantic information, and large sets of grammar rules carefully designed by skilled linguists. This is a very long and therefore costly process.

In the late 1980s when computational power increased and became cheaper, there was more interest in statistical models for machine translation. Statistical models are derived from analysing bilingual text corpora, such as the Europarl parallel corpus, which contains the proceedings of the European Parliament in 11 European languages. Given enough data, statistical MT works well enough to derive an approximate meaning of a foreign language text by processing parallel versions and finding plausible patterns of words. But unlike knowledge-driven systems, statistical (or data-driven) MT often generates ungrammatical output. Data-driven MT is advantageous because less human effort is required, and it can also cover special particularities of the language (e.g., idiomatic expressions) that can get ignored in knowledge-driven systems. 

The strengths and weaknesses of knowledge-driven and data-driven machine translation tend to be complementary, so that nowadays researchers focus on hybrid approaches that combine both methodologies. One approach uses both knowledge-driven and data-driven systems together with a selection module that decides on the best output for each sentence. However, results for sentences longer than say 12 words will often be far from perfect. A better solution is to combine the best parts of each sentence from multiple outputs; this can be fairly complex, as corresponding parts of multiple alternatives are not always obvious and need to be aligned. 

For Bulgarian, MT is particularly challenging. The lack of noun case inflection; free word order and subject pro-drop pose problems for analysis. Extensive inflection in verb morphology is a challenge for generating words with proper markings. 

One of the good examples is WebTrance by SkyCode – a machine translation system which automatically translates texts, help files, menus, windows and internet pages from English, German, French, Spanish, Italian and Turkish into and from Bulgarian.  The use of machine translation can significantly increase productivity provided the system is intelligently adapted to user-specific terminology and integrated into a workflow. 

There is still a huge potential for improving the quality of MT systems. The challenges involve adapting language resources to a given subject domain or user area, and integrating the technology into workflows that already have term bases and translation memories. Another problem is that most of the current systems are English-centred and only support a few languages from and into Bulgarian. This leads to friction in the translation workflow and forces MT users to learn different lexicon coding tools for different systems.

Evaluation campaigns help compare the quality of MT systems, the different approaches and the status of the systems for different language pairs. The table below, which was prepared during the EC Euromatrix+ project, shows the pair-wise performances obtained for 22 of the 23 official EU languages (Irish was not compared). The results are ranked according to a BLEU score, which indicates higher scores for better translations \cite{bleu1}.  A human translator would achieve a score of around 80 points.

The best results (in green and blue) were achieved by languages that benefit from a considerable research effort in coordinated programs and from the existence of many parallel corpora (e.g., English, French, Dutch, Spanish and German). The languages with poorer results are shown in red. These languages either lack such development efforts or are structurally very different from other languages (e.g., Hungarian, Maltese and Finnish).

\subsection{Other Application Areas}

Building language technology applications involves a range of subtasks that do not always surface at the level of interaction with the user, but they provide significant service functionalities "under the hood" of the system in question. They all form important research issues that have now evolved into individual sub-disciplines of computational linguistics. 

Question answering, for example, is an active area of research for which annotated corpora have been built and scientific competitions have been initiated. The concept of question answering goes beyond keyword-based searches (in which the search engine responds by delivering a collection of potentially relevant documents) and enables users to ask a concrete question to which the system provides a single answer. For example:

\hspace{0.5cm}Question: {\it How old was Neil Armstrong when he stepped on the moon?}

\hspace{0.5cm}Answer: {\it 38.}

While question answering is obviously related to the core area of web search, it is nowadays an umbrella term for such research issues as what different types of questions there are, and how they should be handled; how a set of documents that potentially contain the answer can be analysed and compared (do they provide conflicting answers?); and how specific information (the answer) can be reliably extracted from a document without ignoring the context. 

This is in turn related to information extraction (IE), an area that was extremely popular and influential when computational linguistics took a statistical turn in the early 1990s. IE aims to identify specific pieces of information in specific classes of documents, such as detecting the key players in company takeovers as reported in newspaper stories. Another common scenario that has been studied is reports on terrorist incidents. The problem here is to map the text to a template that specifies the perpetrator, target, time, location and results of the incident. Domain-specific template-filling is the central characteristic of IE, which makes it another example of a "behind the scenes" technology that forms a well-demarcated research area that in practice needs to be embedded into a suitable application environment. 

Text summarisation and text generation are two borderline areas that can act either as standalone applications or play a supporting role "under the hood". Summarisation attempts to give the essentials of a long text in a short form, and is one of the features available in Microsoft Word. It mostly uses a statistical approach to identify the "important" words in a text (i.e., words that occur very frequently in the text in question but less frequently in general language use) and determine which sentences contain the most of these "important" words. These sentences are then extracted and put together to create the summary. In this very common commercial scenario, summarisation is simply a form of sentence extraction, and the text is reduced to a subset of its sentences. An alternative approach, for which some research has been carried out, is to generate brand new sentences that do not exist in the source text. This requires a deeper understanding of the text, which means that
  so far this approach is far less robust. On the whole, a text generator is rarely used as a stand-alone application but is embedded into a larger software environment, such as a clinical information system that collects, stores and processes patient data. Creating reports is just one of many applications for text summarisation. 

For the Bulgarian language, research in these text technologies is much less developed than for the English language. Question answering, information extraction, and summarisation have been the focus of numerous open competitions in the USA since the 1990s, primarily organised by the government-sponsored organisations DARPA and NIST. These competitions have significantly improved the start-of-the-art, but their focus has mostly been on the English language. As a result, there are hardly any annotated corpora or other special resources needed to perform these tasks in Bulgarian. When summarisation systems use purely statistical methods, they are largely language-independent and a number of research prototypes are available. For text generation, reusable components have traditionally been limited to surface realisation modules (generation grammars) and most of the available software is for the English language. 

\subsection{Educational Programmes}

Language technology is a very interdisciplinary field that involves the combined expertise of linguists, computer scientists, mathematicians, philosophers, psycholinguists, and neuroscientists among others. As a result, it has not acquired a clear, independent existence in the Bulgarian faculty system. In Bulgarian universities courses in computational linguistics are only partially available. They are usually designed either for humanities students or mathematicians but not both. Two years ago the University of Plovdiv began to offer a Bachelors programme in Linguistics with Information Technologies. 30 students enrolled in 2009-2010 and twice that number in 2010-2011. Students are offered a wide range of courses connected with the fundamentals of linguistics, mathematics and programming as well as language technology applications. The Bachelors degree in Informatics at the University of Plovdiv traditionally offers a lecture course in Computational Linguistics. 

Since 2004 the Faculty of Mathematics and Informatics of the University of Sofia has been offering a Masters programme in Artificial Intelligence. The University of Sofia also offers a Masters programme in Computational Linguistics. The programme includes subjects from the sphere of mathematics, logic, programming and theoretical linguistics. Graduates possess a good grounding to begin academic research in the area of computational linguistics as well as broader computational and linguistic literacy allowing them ability to develop unconventional practical applications. The success of the programme is evident from the professional development of the students from the previous intakes. They are employed in leading IT companies, popular national and specialised media and particularly in academic research.
The education of language technology in sufficient numbers is a prerequisite for the diverse research and thus the development of successful commercial activity.

\subsection{National Projects and Efforts}

The first international and national funding supporting language technologies for Bulgarian began at the very beginning of 1990s. Over a short period of time financing for a number of research projects from European institutions was got: LaTeSLav (Language Processing Technologies for Slavic Languages, 1991-1994) – aimed at developing a prototype of a grammar checker, BILEDITA (Bilingual Electronic Dictionaries and Intelligent Text Alignment, 1996-1998) – for the development of bi-lingual electronic dictionaries, GLOSSER (Support of Second Language Acquisition and Learning from Aligned Corpora, 1996-1998) – aimed at supporting foreign language training and others. The Multext-East (Multilingual Text Tools and Corpora for Central and Eastern European Languages, 1995-1997) extension of the previous Multext and EAGLES (European Commission's Expert Advisory Group on Language Engineering Standards) EU projects provided the Bulgarian language resources in a standardised format
  with standard mark-up and annotation, and these resources were later expanded and upgraded in the ELAN (European Language Activity Network, 1998-1999), TELRI I in II (Trans European Language Resources Infrastructure, 1995-1998 / 1999-2001) and Concede (Consortium for Central European Dictionary Encoding, 1998-2000) projects.

The Ministry of Education, Youth and Research through the National Scientific Fund (NSF), has supported research through national research programs. These programs have impelled numerous research projects and collaboration with international research centres and companies. The basis of technology development and commercial applications for automated processing of the Bulgarian language has been partly created as a result of these projects.

A number of years ago five Bulgarian academic institutions founded a consortium to create and develop an integrated national academic infrastructure for language resources. Bulgarian institutions are also involved in the CLARIN project (Common Language Resources and Technology Infrastructure). Other ongoing projects include those comprised by EUROPEANA aimed at developing the basic technologies and standards necessary to make knowledge on the Internet more widely available in the future. In addition to many other smaller-scale funded projects, the above-mentioned projects have led to the development of competences in the field of Language Technology as well as a basic technological infrastructure of language tools and resources for Bulgarian. However, public funding for LT projects in Bulgaria is dramatically lower than that for comparable projects in Europe, as well as in comparison to investments into areas such as language translation and multilingual information access by
  the USA \cite{sprachtech}.  

\subsection{Availability of Tools and Resources}

The following table summarises the current state of language technology support for the Bulgarian language. The rating for existing tools and resources was generated by leading experts in the field who provided estimates based on a scale from 0 (very low) to 6 (very high) according to seven criteria.

\begin{figure*}[htb]
\centering
%\begin{tabular}{>{\columncolor{orange1}}p{.33\linewidth}ccccccc} % ORIGINAL
\begin{tabular}{>{\columncolor{orange1}}p{.33\linewidth}@{\hspace*{6mm}}c@{\hspace*{6mm}}c@{\hspace*{6mm}}c@{\hspace*{6mm}}c@{\hspace*{6mm}}c@{\hspace*{6mm}}c@{\hspace*{6mm}}c}
\rowcolor{orange1}
 \cellcolor{white}&\begin{sideways}\makecell[l]{Quantity}\end{sideways}
&\begin{sideways}\makecell[l]{\makecell[l]{Availability} }\end{sideways} &\begin{sideways}\makecell[l]{Quality}\end{sideways}
&\begin{sideways}\makecell[l]{Coverage}\end{sideways} &\begin{sideways}\makecell[l]{Maturity}\end{sideways} &\begin{sideways}\makecell[l]{Sustainability}\end{sideways} &\begin{sideways}\makecell[l]{Adaptability}\end{sideways} \\ \addlinespace
\multicolumn{8}{>{\columncolor{orange2}}l}{Language Technology: Tools, Technologies and Applications} \\ \addlinespace
Speech Recognition	&	1.6 &	0.8 &	2.4 &	2.4 &	1.6 &	1.6 &	0.8 \\ \addlinespace
Speech Synthesis &	1.6 &	0.8 &	2.4 &	2.4 &	1.6 &	1.6 &	0.8\\ \addlinespace
Grammatical analysis &	2.4 &	2 &	3.6 &	3.6 &	2.8 &	2.4 &	2.8\\ \addlinespace
Semantic analysis &	0.8 &	0.8 &	1.3 &	1.1 &	0.8 &	1.1 &	1.3\\ \addlinespace
Text generation &	0.8 &	0.8 &	1.6 &	1.6 &	1.6 &	0.8 &	0.8\\ \addlinespace
Machine translation &	2.4 &	1.6 &	1.6 &	1.6 &	1.6 &	1.6 &	2.4\\ \addlinespace
\multicolumn{8}{>{\columncolor{orange2}}l}{Language Resources: Resources, Data and Knowledge Bases} \\ \addlinespace
Text corpora &	2.8 &	2.4 &	3.2 &	2.4 &	3.2 &	2.4 &	2.8\\ \addlinespace
Speech corpora &	0.8 &	0.8 &	2.4 &	1.6 &	2.4 &	2.4 &	2.5\\ \addlinespace
Parallel corpora &	2.4 &	0.8 &	3.2 &	1.6 &	1.6 &	1.6 &	2.4\\ \addlinespace
Lexical resources &	2.4 &	2.8 &	3.6 &	2.8 &	3.2 &	3.2 &	3.2\\ \addlinespace
Grammars &	1.6 &	1.6 &	2.4 &	2.4 &	2.4 &	2.4 &	1.6\\
\end{tabular}
\caption{State of language technology support for Bulgarian}
\label{fig:lrlttable_en}
\end{figure*}

The key results for the Bulgarian language can be summed up as follows:

\begin{itemize}
\item The more linguistic and semantic knowledge a tool draws on, the more gaps there are in the technology (information retrieval versus text semantics). There is a need for far more effort to support deep linguistic processing.

\item Research has been successful in designing particularly high quality software, but many of the resources are not standardised and cannot be sustained effectively. A concerted program is required to standardise data and interchange formats.

\item While some reference corpora of high quality and quantity exist, i.e. Bulgarian national corpus, large syntactically and semantically corpora annotated by experts are not available.  Again, the situation gets worse as the need for more deep linguistic and semantic information grows.

\item There is a lack of parallel corpora for machine translation. Translation from Bulgarian to English works best because this language pair has the most data available. 

\item There is a huge gap in multimedia data.
\end{itemize}

\subsection{Cross-language comparison}

The languages were clustered using the following five-point scale: 

    \begin{itemize}
      \item Cluster 1 (excellent LT support)
      \item Cluster 2 (good support)
      \item Cluster 3 (moderate support)
      \item Cluster 4 (fragmentary support) 
      \item Cluster 5 (weak or no support)
    \end{itemize}

LT support was measured according to the following criteria:

\begin{itemize}
\item Speech Processing: Quality of existing speech recognition technologies, quality of existing speech synthesis technologies, coverage of domains, number and size of existing speech corpora, amount and variety of available speech-based applications
\item Machine Translation: Quality of existing MT technologies, number of language pairs covered, coverage of linguistic phenomena and domains, quality and size of existing parallel corpora, amount and variety of available MT applications
\item Text Analysis: Quality and coverage of existing text analysis technologies (morphology, syntax, semantics), coverage of linguistic phenomena and domains, amount and variety of available applications, quality and size of existing (annotated) text corpora, quality and coverage of existing lexical resources (e.g., WordNet) and grammars
\item Resources: Quality and size of existing text corpora, speech corpora and parallel corpora, quality and coverage of existing lexical resources and grammars
\end{itemize} 

\subsection{Conclusions}

In this series of white papers, we have made an important
initial effort to assess language technology support for 30 European
languages, and provide a high-level comparison across these
languages. By identifying the gaps, needs and deficits, the European
language technology community and related stakeholders are now in a
position to design a large scale research and development programme
aimed at building a truly multilingual, technology-enabled Europe.

We have seen that there are huge differences between Europe’s
languages. While there are good quality software and resources
available for some languages and application areas, others (usually
‘smaller’ languages) have substantial gaps.  Many languages lack basic
technologies for text analysis and the essential resources for
developing these technologies. Others have basic tools and resources
but are as yet unable to invest in semantic processing. We therefore
still need to make a large-scale effort to attain the ambitious goal
of providing high-quality machine translation between all European
languages.

Over the past decade a number of important electronic language resources for Bulgarian (dictionaries, corpora, lexical data bases) as well as programmes for their processing (word sense disambiguation tool, spell checking, etc..) have been developed. However, the scope of the resources and the range of tools are still very limited when compared to the resources and tools for the English language, and they are simply not sufficient in quality and quantity to develop the kind of technologies required to support a truly multilingual knowledge society.

Nor can we simply transfer technologies already developed and optimised for the English language to handle Bulgarian. English-based systems for parsing (syntactic and grammatical analysis of sentence structure) typically perform far less well on German texts, due to the specific characteristics of the Bulgarian language such as free word order or sibject ommision.
The Bulgarian language technology industry dedicated to transforming research into products is currently fragmented and disorganised. A number of specialised small and middle SMEs that are not robust enough to address the internal and the global market with a sustained strategy are working in the field.

  Our findings show that the only alternative is to make a substantial effort to create LT resources for Bulgarian, and use them to drive forward research, innovation and development. The need for large amounts of data and the extreme complexity of language technology systems makes it vital to develop a new infrastructure and a more coherent research organization to spur greater sharing and cooperation.

There is also a lack of continuity in research and development funding. Short-term coordinated programmes tend to alternate with periods of sparse or zero funding. In addition, there is an overall lack of coordination with programmes in other EU countries and at the European Commission level.
In general, it can be stated that in the last two decades language technology for Bulgarian was never supported by a consistently devised national funding scheme. The process of development of HLT applications, tools and resources for Bulgarian has been, therefore, a mixture of international projects extending their scope from Western European languages to Middle and Eastern Europe, also with a view of the EU enlargement process, national research funding, and the enthusiasm of researchers involved in LT.

We can therefore conclude that there is a desperate need for a large, coordinated initiative focused on overcoming the differences in language technology readiness for European languages as a whole. 

Bulgaria’s participation in
META-NET will make it possible to develop, standardise and make
available several important LT resources and thus contribute to the
growth of Bulgarian language technology.

META-NET’s long-term goal is to introduce high-quality language technology for all languages in order to achieve political and economic unity through cultural diversity. The technology will help tear down existing barriers and build bridges between Europe’s languages. This requires all stakeholders — in politics, research, business, and society — to unite their efforts for the future.
\end{multicols}

\clearpage

\ssection[About META-NET]{About META-NET}

\begin{multicols}{2}
                                                                                                                                                                                                                                                                                                                                                META-NET is a Network of Excellence funded by the European Commission. The network currently consists of 47 members from 31 European countries. META-NET fosters the Multilingual Europe Technology Alliance (META), a growing community of language technology professionals and organisations in Europe. 

META-NET cooperates with other initiatives like the Common Language Resources and Technology Infrastructure (CLARIN), which is helping establish digital humanities research in Europe. META-NET fosters the technological foundations for a truly multilingual European information society that:

\begin{itemize}
\item makes communication and cooperation possible across languages;
\item provides equal access to information and knowledge in any language;
\item offers advanced and affordable networked information technology to European citizens.
\end{itemize}

META-NET stimulates and promotes multilingual technologies for all European languages. The technologies enable automatic translation, content production, information processing and knowledge management for a wide variety of applications and subject domains. The network wants to improve current approaches, so better communication and cooperation across languages can take place. Europeans have an equal right to information and knowledge regardless of language. 

META-NET launched on 1 February 2010 with the goal of advancing research in language technology. The network supports a Europe that unites as a single digital market and information space. META-NET has conducted several activities that further its goals. META-VISION, META-SHARE and META-RESEARCH are the network’s three lines of action.

META-VISION fosters a dynamic and influential stakeholder community that unites around a shared vision and a common strategic research agenda (SRA). The main focus of this activity is to build a coherent and cohesive LT community in Europe by bringing together representatives from highly fragmented and diverse groups of stakeholders. In the first year of META-NET, presentations at the FLaReNet Forum (Spain), Language Technology Days (Luxembourg), JIAMCATT 2010 (Luxembourg), LREC 2010 (Malta), EAMT 2010 (France) and ICT 2010 (Belgium) centred on public outreach. According to initial estimates, META-NET has already contacted more than 2,500 LT professionals to develop its goals and visions with them. At the META-FORUM 2010 event in Brussels, META-NET communicated the initial results of its vision building process to more than 250 participants. In a series of interactive sessions, the participants provided feedback on the visions presented by the network. 

META-SHARE creates an open, distributed facility for exchanging and sharing resources. The peer-to-peer network of repositories will contain language data, tools and web services that are documented with high-quality metadata and organised in standardised categories. The resources can be readily accessed and uniformly searched. The available resources include free, open source materials as well as restricted, commercially available, fee-based items. META-SHARE targets existing language data, tools and systems as well as new and emerging products that are required for building and evaluating new technologies, products and services. The reuse, combination, repurposing and re-engineering of language data and tools plays a crucial role. META-SHARE will eventually become a critical part of the LT marketplace for developers, localisation experts, researchers, translators and language professionals from small, mid-sized and large enterprises. META-SHARE addresses the full developmen
 t cycle of LT—from research to innovative products and services. A key aspect of this activity is establishing META-SHARE as an important and valuable part of a European and global infrastructure for the LT community. 

META-RESEARCH builds bridges to related technology fields. This activity seeks to leverage advances in other fields and to capitalise on innovative research that can benefit language technology. In particular, this activity wants to bring more semantics into machine translation (MT), optimise the division of labour in hybrid MT, exploit context when computing automatic translations and prepare an empirical base for MT. META-RESEARCH is working with other fields and disciplines, such as machine learning and the Semantic Web community. META-RESEARCH focuses on collecting data, preparing data sets and organising language resources for evaluation purposes; compiling inventories of tools and methods; and organising workshops and training events for members of the community. This activity has already clearly identified aspects of MT where semantics can impact current best practices. In addition, the activity has created recommendations on how to approach the problem of integrating 
 semantic information in MT. META-RESEARCH is also finalising a new language resource for MT, the Annotated Hybrid Sample MT Corpus, which provides data for English-German, English-Spanish and English-Czech language pairs. META-RESEARCH has also developed software that collects multilingual corpora that are hidden on the Web.
\end{multicols}
